% !TEX root = ../main.tex
\section{Results}

\subsection{LLM-Assisted Screening Performance}

\metagent{} was evaluated on a benchmark comprising 22 published COVID-19 and mpox source-review screening tasks, covering 12,758 candidate records and 668 ground-truth included studies. The benchmark spans three key epidemiological parameters: reproduction number ($n=245$ ground-truth included studies), serial interval ($n=216$), and fatality ($n=207$), which correspond to transmission potential, transmission speed, and disease severity, respectively. The included source reviews and their corresponding raw-pool and ground-truth counts are detailed in Table~\ref{tab:source-reviews}. Candidate records were retrieved from PubMed, and the final inclusion lists reported in the original systematic reviews were used as task-specific ground truth. \metagent{} screens each record from its title and abstract and assigns it to one of three relevance categories: Strong, Possible, or Unlikely. For evaluation, Strong and Possible records were treated as relevant and advanced to manual review.

Across all tasks, \metagent{} achieved an overall recall of 0.960 (95\% CI: 0.942--0.972) while reducing the overall \nns{} from 19.10 in the original candidate pool to 5.25 after LLM screening. Parameter-specific analyses showed consistently high recall, reaching 0.984 (95\% CI: 0.959--0.994) for reproduction number, 0.935 (95\% CI: 0.894--0.961) for serial interval, and 0.957 (95\% CI: 0.919--0.977) for fatality (Figure~\ref{fig:screening-recall-nns}). These results demonstrate that \metagent{} can substantially reduce the number of articles requiring human screening while preserving high sensitivity in systematic review workflows. Model-level comparisons and per-baseline screening metrics, including WR as a secondary operating characteristic, are reported in Figure~\ref{fig:model-compare} and Table~\ref{tab:model-summary}.

The number of FNs was small across all groups. Taking the serial interval group as an example, the 14 FNs primarily originated from records where the target parameter only appeared in the full-text tables or appendices but not in the abstracts. This pattern reflects the inherent information limitation of relying solely on abstracts for the first-stage screening, rather than systematic misjudgment by the model. Thus, it represents an expected and acceptable source of error in the methodological design. In a focused comparison on the Ali et al. serial-interval review, providing full-text passages instead of titles and abstracts increased recall from 0.882 to 0.941 and F1 from 0.662 to 0.780, recovering three previously missed eligible studies and excluding sixteen false positives relative to title-and-abstract screening, indicating that access to full-text evidence can recover false negatives associated with sparse abstract context; experimental setup, detailed counts, and remaining false-positive categories are reported in Appendix C.

Screening efficiency improved across all parameter domains, with the overall \nns{} decreasing from 19.10 to 5.25 and WR reaching 73.6\%. Lower \nns{} indicates fewer records requiring manual review for each GT-positive study retained. Serial interval showed the greatest screening efficiency, reducing \nns{} from 14.11 in the original candidate pool to 2.55 after LLM screening, such that reviewers would screen on average fewer than three records to confirm one eligible study. Reproduction-number screening also reduced \nns{} from 14.13 to 5.70 while achieving the highest recall. Fatality retained high recall and reduced \nns{} substantially from 30.19 to 7.47, although its \nns{} remained higher because eligible fatality studies were less frequent in the candidate pool. Taken together, these results indicate that the residual screening burden is constrained primarily by the density of eligible studies in the candidate pool rather than by the discriminative capacity of the screening model itself.

\begin{figure}[H]
  \centering
  \includegraphics[width=\textwidth]{pdf/screening_recall_nns_overall.pdf}
  \caption{LLM-assisted screening recall and \nns{} (papers needed to screen per included paper) across COVID-19 and mpox benchmarks. a) and c) per-review recall grouped by parameter domain (colour-coded); filled circles are per-review recall; horizontal bars are Wilson 95\% CIs. b) and d) corresponding \nns{}; coloured bars = LLM-assisted, grey bars = original unassisted; whiskers are 95\% CIs.}
  \label{fig:screening-recall-nns}
\end{figure}

\subsection{Full-Text Parameter Coding and Extraction}

Full-text parameter coding and extraction recovered the source-review estimates with close agreement across the benchmark. This component-level analysis used the source-review included studies as fixed inputs, isolating the coding workflow from upstream screening errors, and tested whether localization and codebook-guided extraction could recover the evidence required to reproduce the review-level epidemiological estimates.

Across 22 source reviews, the LLM-coded estimates closely matched the source-review references (Figure~\ref{fig:coding-overall}). Agreement was strongest in reviews for which the target estimand and pooling rule were explicitly defined, as observed for the serial-interval and reproduction-number analyses. For COVID-19, LLM-coded serial-interval estimates differed from the source-review estimates by 0.02--0.09 days across all five reviews, and the reproduction-number estimates preserved the same review-level ordering as the original analyses. Fatality estimates were also recovered when the denominator was aligned with that of the source review, including invasive-mechanical-ventilation fatality, hospitalized fatality, and mpox clade- or population-specific CFR.

These observations were further supported by a supplementary comparison against the LEADS coding baseline (Appendix D). Because the target parameters are measured on different scales, agreement was interpreted primarily within parameter domains. The parameter-stratified MAE was 0.190 days for serial interval, 0.165 for reproduction number, and 0.416 pp for fatality; the mixed-unit overall MAE is reported as a descriptive summary in Appendix D. The codebook-guided workflow produced estimates closer to the source-review reference than LEADS in 21 of 22 source reviews, with the only exception being an incubation-period entry in which both methods remained close to the reference. Taken together, these results indicate that codebook-guided full-text coding can align LLM-coded estimates with source-review epidemiological estimates when the target parameter definition and synthesis rule are made explicit. Full per-review comparisons are reported in Figure~\ref{fig:supp-coding-leads} and Table~\ref{tab:coding-compare-leads-detail}.

\begin{figure}[!t]
  \centering
  \includegraphics[width=\textwidth]{pdf/coding_estimates_overall_by_disease.pdf}
  \caption{Agreement between LLM-coded estimates and source-review estimates. Panels a and b show source-review and LLM-coded estimates for serial interval, reproduction number, and fatality in COVID-19 and mpox reviews. Points denote point estimates and horizontal bars denote 95\% confidence intervals. Grey squares indicate estimates reported in the source reviews; coloured circles indicate estimates derived from LLM-coded full-text evidence. Axes are scaled separately by disease and parameter.}
  \label{fig:coding-overall}
\end{figure}
